% =========================================================
% 5. copy() ORM Method
% =========================================================
\section{\texttt{copy()} ORM Method}

\begin{frame}[standout]
  \texttt{copy()} ORM Method
\end{frame}

\begin{frame}[fragile]{\texttt{COPY()} METHOD}
  \begin{itemize}
    \item The \texttt{copy()} method is an ORM method used to \textbf{duplicate} an existing record.
    \item It creates a new record by copying field values from the current record.
    \item The method returns a recordset containing the newly created (duplicated) record.
  \end{itemize}
  \begin{block}{Method Syntax}
    The standard syntax of the \texttt{copy()} method is:
\begin{lstlisting}
def copy(self, default=None):
    return super().copy(default)
\end{lstlisting}
    This method has three important parts:
    \begin{enumerate}
      \item \texttt{copy}
      \item \texttt{self}
      \item \texttt{default}
    \end{enumerate}
  \end{block}
\end{frame}

\begin{frame}[fragile]{BREAKING DOWN THE METHOD SIGNATURE}
  \begin{itemize}
    \item \texttt{copy}: The ORM method responsible for duplicating records.
    \item \texttt{self}: The source recordset to copy from.
          In normal usage, \texttt{self} is usually \textbf{one} record.
    \item \texttt{default}: An optional dictionary of values that override the copied values.
  \end{itemize}
\begin{lstlisting}
default = {
    'name': 'Ahmed Mohamed (Copy)',
}
\end{lstlisting}
  \begin{itemize}
    \item Keys in \texttt{default} replace the values that would otherwise be copied.
    \item Fields not listed in \texttt{default} are copied from the original record
          (unless marked \texttt{copy=False}).
  \end{itemize}
\end{frame}

\begin{frame}[fragile]{WHAT DOES \texttt{COPY()} RETURN?}
  \begin{itemize}
    \item The \texttt{copy()} method returns a new recordset.
  \end{itemize}
\begin{lstlisting}
student = self.env['student.student'].browse(15)
new_student = student.copy()
print(new_student)
\end{lstlisting}
  Output is (example):
\begin{lstlisting}
student.student(28,)
\end{lstlisting}
  \begin{itemize}
    \item A new database row is created.
    \item The original record remains unchanged.
  \end{itemize}
\end{frame}

\begin{frame}[fragile]{USING THE \texttt{DEFAULT} ARGUMENT}
  \begin{itemize}
    \item You can change some values while copying.
  \end{itemize}
\begin{lstlisting}
student = self.env['student.student'].browse(15)
new_student = student.copy({
    'name': 'Amina Mohamed',
    'email': 'amina@example.com',
})
\end{lstlisting}
  \begin{itemize}
    \item Everything else is copied from student 15.
    \item This is useful for ``Duplicate'' buttons in the UI.
  \end{itemize}
\end{frame}

\begin{frame}[fragile]{WHICH FIELDS ARE COPIED?}
  \begin{itemize}
    \item By default, most stored fields are copied.
    \item A field can opt out of copying:
  \end{itemize}
\begin{lstlisting}
email = fields.Char(copy=False)
registration_code = fields.Char(copy=False)
\end{lstlisting}
  \begin{itemize}
    \item Fields with \texttt{copy=False} become empty / default on the new record.
    \item One2many lines may also be copied depending on the field definition.
    \item Unique constraints may force you to override values through \texttt{default}.
  \end{itemize}
\end{frame}

\begin{frame}[fragile]{METHOD CALL AND OVERRIDE PATTERN}
  \textbf{Method Call}
\begin{lstlisting}
new_student = student.copy({'name': 'Ahmed (Copy)'})
\end{lstlisting}
  \textbf{Common Override Pattern}
\begin{lstlisting}
def copy(self, default=None):
    default = dict(default or {})
    default.setdefault('name', f"{self.name} (Copy)")
    return super().copy(default)
\end{lstlisting}
  \begin{itemize}
    \item Overrides are useful when duplicated names must stay unique.
  \end{itemize}
\end{frame}
